\documentclass[11pt]{article}

\usepackage[margin=2.5cm]{geometry}
\usepackage{amsmath,amssymb,amsthm,bm}
\usepackage{enumitem}

\title{Exercise Sheet -- Probability Spaces \& Random Variables\\[0.4em]
{\large (Events, Laws, Expectation, Variance, Inequalities, and Basic Tools)}}
\author{}
\date{}

\begin{document}
\maketitle

This sheet is organized in two parts:
\begin{itemize}
  \item \textbf{Part A -- Fundamental exercises (9 problems)}
  \item \textbf{Part B -- Advanced exercises (11 problems)}
\end{itemize}

The problems are designed to practice:
events as sets, probability measures and their properties (union bound, total probability, Bayes),
random variables and distributions (pmf/pdf/cdf), expectation/variance and transformations,
and classical inequalities (Jensen/Markov/Chebyshev), with a few finance-flavored applications.

\bigskip
\hrule
\bigskip

\section*{Part A -- Fundamental Exercises}

\begin{enumerate}[leftmargin=*, itemsep=1.2em]

% -----------------------------
\item \textbf{Events as sets: operations, unions/intersections, and $\limsup/\liminf$}

\textbf{(a) Dice warm-up.}
Let $\Omega=\{1,2,3,4,5,6\}$ and
\[
A=\{2\},\qquad B=\{2,4,6\},\qquad C=\{1,2,3\}.
\]
Compute $A^c,B^c,C^c$, then $A\cup C$, $B\cap C$, and $B\setminus C$.
Verify that $B\setminus C=B\cap C^c$.

\medskip
\textbf{(b) Infinite unions/intersections and limits.}
For $n\ge1$, define
\[
A_n=\Big[-\frac1n,\;2+\frac1n\Big],
\qquad
B_n=\Big[-\frac5n,\;\frac{n}{2}\Big].
\]
\begin{enumerate}[label=(\roman*)]
\item Compute $\bigcup_{n\ge1}A_n$ and $\bigcap_{n\ge1}A_n$.
\item Compute $\limsup_{n\to\infty}A_n$ and $\liminf_{n\to\infty}A_n$ and interpret them.
\item Same questions for $(B_n)_{n\ge1}$.
\end{enumerate}

\bigskip

% -----------------------------
\item \textbf{Basic probability properties + union bound (proof + application)}

Let $(\Omega,\mathcal{F},\mathbb{P})$ be a probability space and $A,B\in\mathcal{F}$.
\begin{enumerate}[label=(\alph*)]
\item Prove: if $A\subseteq B$ then $\mathbb{P}(A)\le \mathbb{P}(B)$.
\item Prove: $\mathbb{P}(A^c)=1-\mathbb{P}(A)$.
\item Prove inclusion--exclusion:
\[
\mathbb{P}(A\cup B)=\mathbb{P}(A)+\mathbb{P}(B)-\mathbb{P}(A\cap B),
\]
and deduce $\mathbb{P}(A\cup B)\le \mathbb{P}(A)+\mathbb{P}(B)$.
\item (Union bound, finite case) Show that for events $A_1,\dots,A_n$,
\[
\mathbb{P}\Big(\bigcup_{k=1}^n A_k\Big)\le \sum_{k=1}^n \mathbb{P}(A_k).
\]
\end{enumerate}

\bigskip

% -----------------------------
\item \textbf{Independence: pairwise vs mutual (finite example)}

Let $\Omega=\{\omega_1,\omega_2,\omega_3,\omega_4\}$ with the uniform probability.
Define
\[
A=\{\omega_1,\omega_2\},\quad
B=\{\omega_1,\omega_3\},\quad
C=\{\omega_2,\omega_3\}.
\]
\begin{enumerate}[label=(\alph*)]
\item Show that $A,B,C$ are pairwise independent.
\item Compute $\mathbb{P}(A\cap B\cap C)$ and compare with $\mathbb{P}(A)\mathbb{P}(B)\mathbb{P}(C)$.
\item Explain briefly what this teaches you about ``pairwise'' vs ``mutual'' independence.
\end{enumerate}

\bigskip

% -----------------------------
\item \textbf{Independence and dependence (coin tosses)}

Let $(X_1,X_2,X_3)$ be outcomes of three independent fair coin tosses, $X_i\in\{H,T\}$.
Define
\[
A=\{X_1=H\},\qquad
B=\{X_2=H\},\qquad
C=\{\text{exactly two heads occur}\}.
\]
\begin{enumerate}[label=(\alph*)]
\item Show that $A$ and $B$ are independent.
\item Are $A$ and $C$ independent? Justify by computing probabilities.
\item Give an example of two events built from $(X_1,X_2,X_3)$ that are \emph{not} independent.
\end{enumerate}

\bigskip

% -----------------------------
\item \textbf{Conditioning: the two-children problem}

A family has two children. Assume the four configurations
\[
(F,F), (F,G), (G,F), (G,G)
\]
are equiprobable ($F$=girl, $G$=boy; first coordinate = older child).
Compute:
\begin{enumerate}[label=(\alph*)]
\item $\mathbb{P}(\text{both girls}\mid \text{younger is a girl})$,
\item $\mathbb{P}(\text{both girls}\mid \text{older is a girl})$,
\item $\mathbb{P}(\text{both girls}\mid \text{at least one is a girl})$.
\end{enumerate}

\bigskip

% -----------------------------
\item \textbf{Law of total probability + Bayes (finite case)}

A market indicator $Y\in\{G,B\}$ (Good/Bad) satisfies
$\mathbb{P}(Y=G)=0.7$ and $\mathbb{P}(Y=B)=0.3$.
A signal $S\in\{1,0\}$ is observed with
\[
\mathbb{P}(S=1\mid Y=G)=0.9,
\qquad
\mathbb{P}(S=1\mid Y=B)=0.2.
\]
\begin{enumerate}[label=(\alph*)]
\item Compute $\mathbb{P}(S=1)$ using the law of total probability.
\item Compute $\mathbb{P}(Y=G\mid S=1)$ using Bayes' formula.
\item Interpret the result in one sentence.
\end{enumerate}

\bigskip

% -----------------------------
\item \textbf{Bernoulli/binomial modeling (oil slicks)}

In an oil region, the probability that one drilling leads to an oil slick is $0.1$.
\begin{enumerate}[label=(\alph*)]
\item Explain why one drilling can be modeled by a Bernoulli random variable.
\item We perform $10$ drillings. Let $X$ be the number of oil slicks.
Under which assumptions does $X\sim\mathrm{Bin}(10,0.1)$ hold?
\item Assuming $X\sim\mathrm{Bin}(10,0.1)$, compute $\mathbb{P}(X=2)$ and $\mathbb{P}(X\ge1)$.
\end{enumerate}

\bigskip

% -----------------------------
\item \textbf{CDF, expectation, and variance for a discrete law}

Let $X$ satisfy
\[
\mathbb{P}(X=-1)=\frac14,\quad
\mathbb{P}(X=0)=\frac12,\quad
\mathbb{P}(X=2)=\frac14.
\]
\begin{enumerate}[label=(\alph*)]
\item Write the cdf $F_X(t)=\mathbb{P}(X\le t)$ explicitly (piecewise).
\item Check that $F_X$ is non-decreasing and right-continuous, and identify its jumps.
\item Compute $\mathbb{P}(-1<X\le2)$ using the cdf.
\item Compute $\mathbb{E}[X]$ and $\mathrm{Var}(X)$.
\end{enumerate}

\bigskip

% -----------------------------
\item \textbf{Mean/variance for standard laws (short answers)}

Compute mean and variance for:
\begin{enumerate}[label=(\alph*)]
\item $X\sim \mathrm{Bin}(n,p)$,
\item $Y\sim \mathrm{Poi}(\lambda)$,
\item $U\sim U[a,b]$,
\item $E\sim \mathrm{Exp}(\lambda)$,
\item $Z\sim \mathcal{N}(\mu,\sigma^2)$.
\end{enumerate}
(You may state known formulas, but justify at least two of them.)

\bigskip

% -----------------------------
\item \textbf{Transformations}

\textbf{(a) Continuous transformation.}
Let $X\sim U[0,1]$ and define $Z=X^2$.
Compute $F_Z$, deduce $f_Z$, and compute $\mathbb{E}[Z]$ in two ways.

\medskip
\textbf{(b) Min/max.}
Let $X\sim U[0,1]$. Define
\[
Y=\min(X,1-X),\qquad W=\max(X,1-X).
\]
Determine the cdf (or density) of $Y$ and of $W$, and compute $\mathbb{E}[YW]$.

\bigskip

\bigskip

% -----------------------------
\item \textbf{Quadratic transformation of a normal variable}

Let $X\sim \mathcal{N}(0,1)$ and define
\[
Y = X^2.
\]

\begin{enumerate}[label=(\alph*)]
\item Compute the cdf $F_Y(y)=\mathbb{P}(Y\le y)$ for $y\ge0$.
\item Deduce the density $f_Y$.
\item Identify the distribution of $Y$.
\item Compute $\mathbb{E}[Y]$ and $\mathrm{Var}(Y)$.
\end{enumerate}

\bigskip

% -----------------------------
\item \textbf{Absolute value of a normal variable}

Let $X\sim\mathcal{N}(0,\sigma^2)$ and define
\[
Y=|X|.
\]

\begin{enumerate}[label=(\alph*)]
\item Compute the cdf of $Y$.
\item Deduce the density of $Y$.
\item Compute $\mathbb{E}[Y]$.
\item Explain why $Y$ is not Gaussian.
\end{enumerate}

\bigskip

% -----------------------------
\item \textbf{Characteristic function and identification (detailed)}

Let $X$ be a real random variable with characteristic function
\[
\varphi_X(t)=\mathbb{E}\!\left[e^{itX}\right]=\frac{1}{1+t^2},\qquad t\in\mathbb{R}.
\]

\begin{enumerate}[label=(\alph*)]

\item \textbf{Basic checks.}
\begin{enumerate}[label=(\roman*)]
\item Verify that $\varphi_X(0)=1$ and that $\varphi_X$ is even and real-valued.
\item Show that $|\varphi_X(t)|\le 1$ for all $t$.
\item Compute $\displaystyle \lim_{|t|\to\infty}\varphi_X(t)$.
\end{enumerate}

\medskip

\item \textbf{Integrability.}
Show that $\varphi_X\in L^1(\mathbb{R})$ by computing
\[
\int_{-\infty}^{\infty}\left|\varphi_X(t)\right|\,dt
=
\int_{-\infty}^{\infty}\frac{1}{1+t^2}\,dt.
\]
(You may recall $\int_{-\infty}^{\infty}\frac{1}{1+t^2}\,dt=\pi$.)

\medskip

\item \textbf{Recovering a density via Fourier inversion.}
Assume that $\varphi_X\in L^1(\mathbb{R})$. Then $X$ admits a continuous density $f_X$ given by
\[
f_X(x)=\frac{1}{2\pi}\int_{-\infty}^{\infty} e^{-itx}\,\varphi_X(t)\,dt.
\]
\begin{enumerate}[label=(\roman*)]
\item Using the known Fourier transform identity
\[
\int_{-\infty}^{\infty} e^{-itx}\,\frac{1}{1+t^2}\,dt
= \pi e^{-|x|},
\]
deduce that
\[
f_X(x)=\frac{1}{2}e^{-|x|},\qquad x\in\mathbb{R}.
\]
\item Check that $f_X$ integrates to $1$.
\end{enumerate}

\medskip

\item \textbf{Identification of the law.}
\begin{enumerate}[label=(\roman*)]
\item Identify the distribution corresponding to $f_X(x)=\frac12 e^{-|x|}$.
(State its name and parameters.)
\item Compute $\mathbb{P}(X\ge 0)$ and explain why it must be $1/2$.
\end{enumerate}

\medskip

\item \textbf{Moments: existence and computation.}
\begin{enumerate}[label=(\roman*)]
\item Does $\mathbb{E}[X]$ exist? Justify carefully.
\item Compute $\mathbb{E}[X]$.
\item Compute $\mathbb{E}[|X|]$ and $\mathrm{Var}(X)$.
\end{enumerate}

\medskip

\item \textbf{Link between smoothness of $\varphi_X$ and moments (optional but instructive).}
\begin{enumerate}[label=(\roman*)]
\item Compute $\varphi_X'(t)$ and check whether $\varphi_X'(0)$ exists.
Recall that if $\mathbb{E}[|X|]<\infty$, then $\varphi_X'(0)=i\mathbb{E}[X]$.
\item Compute $\varphi_X''(0)$ and discuss what it suggests about $\mathbb{E}[X^2]$.
(Recall: if $\mathbb{E}[X^2]<\infty$ then $\varphi_X''(0)=-\mathbb{E}[X^2]$.)
\end{enumerate}

\medskip

\item \textbf{Comparison with the Gaussian case.}
Let $Z\sim\mathcal{N}(0,\sigma^2)$ so that
\[
\varphi_Z(t)=\exp\!\left(-\frac{\sigma^2 t^2}{2}\right).
\]
\begin{enumerate}[label=(\roman*)]
\item Compare the decay rates of $\varphi_X(t)=1/(1+t^2)$ and $\varphi_Z(t)$ as $|t|\to\infty$.
\item Explain qualitatively what this implies about the tails of $X$ versus $Z$.
\item Which variable has heavier tails? Give one quantitative argument (e.g.\ compare $\mathbb{P}(|X|>x)$ and $\mathbb{P}(|Z|>x)$ asymptotically).
\end{enumerate}

\end{enumerate}
\bigskip
% -----------------------------
\item \textbf{Non-monotone transformation}

Let $X\sim U[-1,2]$ and define
\[
Y = X(1-X).
\]

\begin{enumerate}[label=(\alph*)]
\item Determine the range of $Y$.
\item Compute $F_Y(y)=\mathbb{P}(Y\le y)$.
\item Deduce the density of $Y$.
\item Sketch the density.
\end{enumerate}

\bigskip

% -----------------------------
\item \textbf{Exponential transformation (finance flavor)}

Let $X\sim \mathcal{N}(\mu,\sigma^2)$ and define
\[
S = e^X.
\]

\begin{enumerate}[label=(\alph*)]
\item Compute the cdf of $S$.
\item Deduce the density of $S$.
\item Compute $\mathbb{E}[S]$.
\item Identify the distribution of $S$ and interpret it financially.
\end{enumerate}

\bigskip

% -----------------------------
\item \textbf{Minimum of two independent variables}

Let $X,Y$ be independent $\mathrm{Exp}(\lambda)$ random variables and define
\[
Z=\min(X,Y).
\]

\begin{enumerate}[label=(\alph*)]
\item Compute $\mathbb{P}(Z>t)$.
\item Deduce the distribution of $Z$.
\item Compute $\mathbb{E}[Z]$.
\item Give a short interpretation (waiting times).
\end{enumerate}

\bigskip

% -----------------------------
\item \textbf{Maximum of two independent variables}

Let $X,Y$ be independent $U[0,1]$ variables and define
\[
M=\max(X,Y).
\]

\begin{enumerate}[label=(\alph*)]
\item Compute the cdf of $M$.
\item Deduce the density of $M$.
\item Compute $\mathbb{E}[M]$.
\item Compare $\mathbb{E}[M]$ with $\mathbb{E}[X]$.
\end{enumerate}

\bigskip

% -----------------------------
\item \textbf{Ratio of uniforms}

Let $X,Y$ be independent $U[0,1]$ random variables and define
\[
Z=\frac{X}{Y}.
\]

\begin{enumerate}[label=(\alph*)]
\item Compute $F_Z(z)$ for $z\ge0$.
\item Deduce the density of $Z$.
\item Does $\mathbb{E}[Z]$ exist? Justify.
\item Explain briefly why heavy tails appear.
\end{enumerate}

\bigskip

% -----------------------------
\item \textbf{Characteristic function and identification}

Let $X$ be a random variable with characteristic function
\[
\varphi_X(t)=\frac{1}{1+t^2}.
\]

\begin{enumerate}[label=(\alph*)]
\item Show that $\varphi_X$ is integrable.
\item Identify the distribution of $X$.
\item Does $\mathbb{E}[X]$ exist?
\item Compare with the Gaussian case.
\end{enumerate}

\bigskip

% -----------------------------
\item \textbf{Censored payoff (option-style)}

Let $X\sim \mathcal{N}(0,1)$ and define
\[
Y=\max(X,0).
\]

\begin{enumerate}[label=(\alph*)]
\item Compute the cdf of $Y$.
\item Show that $Y$ has a mixed distribution.
\item Compute $\mathbb{E}[Y]$.
\item Interpret $Y$ as a financial payoff.
\end{enumerate}

\bigskip

% -----------------------------
\item \textbf{Log-return aggregation}

Let $X_1,\dots,X_n$ be i.i.d.\ $\mathcal{N}(\mu,\sigma^2)$ and define
\[
S_n=\sum_{k=1}^n X_k.
\]

\begin{enumerate}[label=(\alph*)]
\item Compute the distribution of $S_n$.
\item Compute $\mathbb{E}[e^{S_n}]$.
\end{enumerate}

\bigskip

% -----------------------------
\item \textbf{Censored exponential: mixed distribution}

Let $X\sim \mathrm{Exp}(1)$ and fix $a>0$. Define $Y=\min(X,a)$.
\begin{enumerate}[label=(\alph*)]
\item Compute the cdf of $Y$ and sketch it.
\item Does $Y$ admit a density on $\mathbb{R}$? Explain clearly (continuous part vs atom).
\item Compute $\mathbb{E}[Y]$ (hint: $Y=X\mathbf{1}_{X\le a}+a\mathbf{1}_{X>a}$).
\end{enumerate}

\bigskip

% -----------------------------
\item \textbf{Markov and Chebyshev: basic bounds}

Let $X\ge0$ a.s.\ with $\mathbb{E}[X]=5$.
Use Markov's inequality to bound $\mathbb{P}(X\ge 20)$.

Now assume $X$ is integrable with $\mathbb{E}[X]=2$ and $\mathrm{Var}(X)=9$.
Use Chebyshev's inequality to bound $\mathbb{P}(|X-2|\ge 6)$.
Give one sentence explaining why these bounds can be loose.

\end{enumerate}

\bigskip
\hrule
\bigskip

\section*{Part B -- Advanced Exercises}

\begin{enumerate}[leftmargin=*, itemsep=1.2em]
\setcounter{enumi}{9} % continue numbering

% -----------------------------
\item \textbf{Limsup/liminf of events and interpretation}

Let $(A_n)_{n\ge 1}$ be events. Recall:
\[
\limsup_{n\to\infty}A_n=\bigcap_{n\ge1}\ \bigcup_{k\ge n} A_k,
\qquad
\liminf_{n\to\infty}A_n=\bigcup_{n\ge1}\ \bigcap_{k\ge n} A_k.
\]
\begin{enumerate}[label=(\alph*)]
\item Show that $\liminf A_n \subseteq \limsup A_n$.
\item Explain in words what the events $\limsup A_n$ and $\liminf A_n$ mean.
\item Give an example of a sequence $(A_n)$ such that $\liminf A_n=\emptyset$ but $\limsup A_n\neq\emptyset$.
\end{enumerate}

% -----------------------------
\item \textbf{Monotone sequences of events (continuity of probability)}

Let $(A_n)$ be increasing ($A_n\subseteq A_{n+1}$) and $(B_n)$ decreasing ($B_n\supseteq B_{n+1}$).
\begin{enumerate}[label=(\alph*)]
\item Prove that $\mathbb{P}(\cup_{n\ge1}A_n)=\lim_{n\to\infty}\mathbb{P}(A_n)$.
\item Prove that $\mathbb{P}(\cap_{n\ge1}B_n)=\lim_{n\to\infty}\mathbb{P}(B_n)$.
\item Example: Let $\Omega=(0,1]$ with $\mathbb{P}$ uniform, and $A_n=(0,1/n]$.
Compute $\mathbb{P}(A_n)$ and $\mathbb{P}(\cap_{n\ge1}A_n)$, and check consistency with (b).
\end{enumerate}

% -----------------------------
\item \textbf{Borel--Cantelli (I) and a concrete example}

Let $(A_n)_{n\ge1}$ be events such that $\sum_{n=1}^\infty \mathbb{P}(A_n)<\infty$.
\begin{enumerate}[label=(\alph*)]
\item State precisely the conclusion of Borel--Cantelli (I) in terms of $\limsup A_n$.
\item Example: let $(X_n)_{n\ge1}$ be i.i.d.\ with $\mathbb{P}(X_n=1)=1/n^2$ and $\mathbb{P}(X_n=0)=1-1/n^2$.
Let $A_n=\{X_n=1\}$. Show that $\mathbb{P}(A_n\ \text{i.o.})=0$.
\item Interpret the result informally: what does it say about the number of times $X_n=1$ occurs?
\end{enumerate}

% -----------------------------
\item \textbf{Borel--Cantelli (II) under independence}

Assume $(A_n)_{n\ge1}$ are \emph{independent} events and $\sum_{n=1}^\infty \mathbb{P}(A_n)=\infty$.
\begin{enumerate}[label=(\alph*)]
\item State Borel--Cantelli (II).
\item Example: independent coin tosses $(X_n)$ with $\mathbb{P}(X_n=H)=1/2$ and $A_n=\{X_n=H\}$.
Use (a) to conclude $\mathbb{P}(A_n\ \text{i.o.})=1$.
\item Now let $A_n=\{X_n=H \ \text{and}\ X_{n+1}=H\}$. Are the events $(A_n)$ independent?
Can you still argue that $A_n$ happens infinitely often with probability $1$? (Give a short justification.)
\end{enumerate}


% -----------------------------
\item \textbf{Expectation may not exist (Cauchy)}

Let $X$ have the Cauchy density $f(x)=\frac{1}{\pi(1+x^2)}$.
\begin{enumerate}[label=(\alph*)]
\item Show that $\int_{\mathbb{R}} |x| f(x)\,dx = \infty$.
\item Deduce that $\mathbb{E}[X]$ does not exist (in the Lebesgue sense).
\item (Optional) Explain why symmetry alone does \emph{not} guarantee existence of $\mathbb{E}[X]$.
\end{enumerate}

% -----------------------------
\item \textbf{Jensen: a useful bound}

Let $X\ge 0$ be integrable.
\begin{enumerate}[label=(\alph*)]
\item Apply Jensen's inequality to $\varphi(x)=x^2$ and deduce $\mathbb{E}[X]^2\le \mathbb{E}[X^2]$.
\item Apply Jensen to $\varphi(x)=e^x$ and deduce $e^{\mathbb{E}[X]}\le \mathbb{E}[e^X]$.
\item Give a one-line interpretation of (b) in terms of ``averaging then exponentiating'' vs.\ ``exponentiating then averaging''.
\end{enumerate}

% -----------------------------
\item \textbf{Joint density, marginals, and independence check}

Let $(X,Y)$ have joint density
\[
f_{X,Y}(x,y)=6x^2y\,\mathbf{1}_{\{0\le x\le 1,\ 0\le y\le 1\}}.
\]
\begin{enumerate}[label=(\alph*)]
\item Check that $f_{X,Y}$ integrates to $1$ on $[0,1]^2$.
\item Compute $f_X(x)$ and $f_Y(y)$.
\item Show that $X$ and $Y$ are independent.
\item Compute $\mathbb{E}[X]$, $\mathbb{E}[Y]$, and $\mathrm{Var}(X)$.
\end{enumerate}

% -----------------------------
\item \textbf{Covariance and correlation basics}

Let $X,Y\in L^2$. Recall
\[
\mathrm{Cov}(X,Y)=\mathbb{E}[(X-\mathbb{E}X)(Y-\mathbb{E}Y)].
\]
\begin{enumerate}[label=(\alph*)]
\item Show that $\mathrm{Var}(X+Y)=\mathrm{Var}(X)+\mathrm{Var}(Y)+2\,\mathrm{Cov}(X,Y)$.
\item Deduce that if $X$ and $Y$ are independent and square-integrable, then
$\mathrm{Var}(X+Y)=\mathrm{Var}(X)+\mathrm{Var}(Y)$.
\item Give a counterexample where $\mathrm{Cov}(X,Y)=0$ but $X$ and $Y$ are not independent.
(Hint: take $X$ symmetric and $Y=X^2$.)
\end{enumerate}

% -----------------------------
\item \textbf{Finance: portfolio mean and variance}

Let $R=(R_1,\dots,R_d)$ be a vector of asset returns with mean $\mu=\mathbb{E}[R]$
and covariance matrix $\Sigma=\mathrm{Cov}(R)$.
For weights $w\in\mathbb{R}^d$, define the portfolio return $R_p=w^\top R$.
\begin{enumerate}[label=(\alph*)]
\item Show that $\mathbb{E}[R_p]=w^\top \mu$.
\item Show that $\mathrm{Var}(R_p)=w^\top\Sigma w$.
\item Two-asset case: $d=2$, $\Sigma=\begin{pmatrix}\sigma_1^2&\rho\sigma_1\sigma_2\\ \rho\sigma_1\sigma_2&\sigma_2^2\end{pmatrix}$,
and $w_2=1-w_1$. Express $\mathrm{Var}(R_p)$ as a quadratic function of $w_1$ and find the minimizer $w_1^\star$.
\end{enumerate}

% -----------------------------
\item \textbf{Finding a distribution via cdf (practice)}

Let $X$ be continuous with density $f_X(x)=2x\,\mathbf{1}_{[0,1]}(x)$.
Define $Y=\sqrt{X}$.
\begin{enumerate}[label=(\alph*)]
\item Compute $F_X(t)$ for all $t$.
\item Compute $F_Y(t)$ for all $t$ by writing $\mathbb{P}(Y\le t)$ in terms of $X$.
\item Deduce the density $f_Y$.
\item Compute $\mathbb{E}[Y]$.
\end{enumerate}

\end{enumerate}

\end{document}